\section*{Capítulo \ref{ch:b3:palaeomagnetism} --- Geomagnetismo y paleomagnetismo}

\begin{solution}{exo:b3:palaeomagnetism:1}
$\tan I = 2\tan\ang{30} = 1.1547$, por tanto $I = \ang{49.1}$. Recíprocamente, $\tan\lambda
= \tfrac12\tan(\ang{-20}) = -0.1820$, por tanto $\lambda = \ang{-10.3}$, es decir,
\ang{10.3}\,S.
\end{solution}

\begin{solution}{exo:b3:palaeomagnetism:2}
$H = \sqrt{20000^2 + 2000^2} = \qty{20100}{nT}$;
$F = \sqrt{20000^2 + 2000^2 + 40000^2} = \qty{44770}{nT}$. La declinación magnética es
$\arctan(-2000/20000) = \ang{-5.7}$, es decir, \ang{5.7} al oeste del norte;
la inclinación es $\arctan(40000/20100) = \ang{63.3}$, hacia abajo.
\end{solution}

\begin{solution}{exo:b3:palaeomagnetism:3}
$B_0 = -g_1^0 = \qty{29350}{nT}$ en el ecuador, y $2B_0 = \qty{58700}{nT}$ en
los polos. $m = 4\pi a^3B_0/\mu_0 = a^3B_0/\num{e-7} = \num{2.5862e20}\times
\num{2.935e-5}/\num{e-7} = \qty{7.59e22}{A.m^2}$, algo menos que
\qty{7.69e22}{A.m^2} del dipolo completo, que añade $g_1^1$ y $h_1^1$.
\end{solution}

\begin{solution}{exo:b3:palaeomagnetism:4}
(a) Una \omterm{def:b3:palaeomagnetism:nrm}{magnetización termorremanente}, que registra el campo a medida que la lava se enfrió
a través de las \omterm{def:b3:palaeomagnetism:blocking}{temperaturas de bloqueo} de sus granos. (b) Una \omterm{def:b3:palaeomagnetism:nrm}{magnetización remanente
detrítica}, que registra el campo durante el tiempo que tardó el limo en depositarse y
quedar fijado por compactación. (c) Una \omterm{def:b3:palaeomagnetism:nrm}{magnetización remanente química}, que registra el
campo durante el crecimiento de la hematites, durante el enterramiento temprano y, por tanto, algo
después de la sedimentación. (d) Una remanencia viscosa en la dirección del campo actual, añadida a la
TRM original por los granos de tiempo de relajación más corto; se elimina con los
primeros pasos de desmagnetización.
\end{solution}

\begin{solution}{exo:b3:palaeomagnetism:5}
Los vectores unitarios $(\cos I\cos D,\ \cos I\sin D,\ \sin I)$ son
\begin{gather*}
  (0.6545, 0.1391, 0.7431),\quad (0.6670, 0.2695, 0.6947),\\
  (0.5235, 0.1501, 0.8387),\quad (0.5985, 0.0629, 0.7986).
\end{gather*}
La suma es $(2.4435, 0.6216, 3.0751)$, de longitud $R = 3.9766$. Media:
$D = \arctan(0.6216/2.4435) = \ang{14.3}$, $I = \arcsin(3.0751/3.9766) =
\ang{50.7}$. $k = 3/(4 - 3.9766) = 128$. Con $20^{1/3} - 1 = 1.7144$,
$\cos\alpha_{95} = 1 - (0.0234/3.9766)\times1.7144 = 0.98991$, por tanto
$\alpha_{95} = \ang{8.1}$: con cuatro especímenes el cono se mantiene amplio, por muy
compacto que sea el grupo.
\end{solution}

\begin{solution}{exo:b3:palaeomagnetism:6}
\sloppy $\lambda = \arctan(\tfrac12\tan\ang{50}) = \ang{30.79}$, por tanto $p = \ang{59.21}$.
$\sin\lambda_p = \sin\ang{40}\cos\ang{59.21} + \cos\ang{40}\sin\ang{59.21}
\cos\ang{20} = 0.9474$, $\lambda_p = \ang{71.34}$\,N. $\sin\beta =
\sin\ang{59.21}\sin\ang{20}/\cos\ang{71.34} = 0.9184$, $\beta = \ang{66.66}$.
Como $\cos p = 0.5119 < \sin\ang{40}\times0.9474 = 0.6090$, se aplica el segundo caso:
$\varphi_p = 10 + 180 - 66.66 = \ang{123.34}$\,E. El polo está en
\ang{71.3}\,N, \ang{123.3}\,E.
\end{solution}

\begin{solution}{exo:b3:palaeomagnetism:7}
(a) $(7.69 - 8.32)/8.32 = \qty{-7.6}{\percent}$ en 125 años, o
\qty{-6.1}{\percent} por siglo. (b) $29350/12.6 = 2329$, unos 2300 años.
No es una predicción: la tasa misma cambia de un modelo de cinco años al
siguiente, el momento ha sido tanto mayor como menor en el pasado sin
invertirse, y el campo se invierte a intervalos irregulares que ninguna tasa
predice.
\end{solution}

\begin{solution}{exo:b3:palaeomagnetism:8}
Entre 20 y \qty{300}{\celsius} el vector perdió $(6.4, 0.0, 13.6)$: una
componente de \qty{15.0}{mA/m}, con $D = \ang{0.0}$ e
$I = \arctan(13.6/6.4) = \ang{64.8}$. Por encima de \qty{300}{\celsius} perdió $(25.8, 18.1, 24.6)$:
\qty{40.0}{mA/m}, $D = \arctan(18.1/25.8) = \ang{35.1}$, $I =
\arctan(24.6/31.5) = \ang{38.0}$. La segunda, que llega al origen, es la
ChRM; la primera, según una dirección hacia el norte y muy inclinada, que se elimina a baja temperatura,
es una superimposición viscosa a lo largo del campo reciente.
\end{solution}

\begin{solution}{exo:b3:palaeomagnetism:9}
$\lambda = \arctan(\tfrac12\tan(\ang{-35})) = \ang{-19.3}$. Para $I = \ang{-40}$
y \ang{-30}: \ang{-22.8} y \ang{-16.1}, un intervalo de \ang{6.7}.
$p = \ang{90} - \lambda = \ang{109.3}$, $\cos^2p = 0.1096$:
$\mathrm dp = 2.5\times(1 + 0.329) = \ang{3.3}$ y $\mathrm dm =
5\times\sin\ang{109.3}/\cos\ang{35} = \ang{5.8}$.
\end{solution}

\begin{solution}{exo:b3:palaeomagnetism:10}
(a) $\tan\lambda = \tfrac12$, $\lambda = \ang{26.6}$. $F$ es $B_0$ en el
ecuador y $2B_0$ en los polos: una razón de 2. (b) $\tan I = 2\tan\lambda$
da $(1 + \tan^2I)\,\mathrm dI = 2(1 + \tan^2\lambda)\,\mathrm d\lambda$, por tanto
\[
  \frac{\mathrm dI}{\mathrm d\lambda} = \frac{2/\cos^2\lambda}{1 + 4\tan^2\lambda}
  = \frac{2}{\cos^2\lambda + 4\sin^2\lambda} = \frac{2}{1 + 3\sin^2\lambda}.
\]
Un error $\delta I$ da $\delta\lambda = \tfrac12(1 + 3\sin^2\lambda)\,\delta I$,
desde $\tfrac12\delta I$ en el ecuador hasta $2\,\delta I$ en los polos: el
mayor error se produce cerca de los polos. (c) Cerca del ecuador la inclinación cambia
el doble de rápido que la latitud, por lo que un error dado de inclinación cuesta la mitad
de latitud; cerca de los polos la inclinación apenas cambia con la latitud.
\end{solution}

\begin{solution}{exo:b3:palaeomagnetism:11}
(a) $\cos p = \sin(\ang{-35})\sin\ang{70} + \cos\ang{35}\cos\ang{70}\cos\ang{150}
= -0.5390 - 0.2426 = -0.7816$, $p = \ang{141.41}$. Entonces $\tan I_x = 2\cot p =
-2.506$, $I_x = \ang{-68.2}$, y $\lambda_x = \ang{90} - p = \ang{-51.4}$.
$\cos D_x = (\sin\ang{70} - \sin(\ang{-35})\cos p)/(\cos\ang{35}\sin p) =
0.9617$, y como el polo está al este del meridiano del yacimiento ($\sin\ang{150} >
0$), $D_x = \ang{15.9}$. (b) $\lambda_o = \arctan(\tfrac12\tan(\ang{-60})) =
\ang{-40.9}$. $R = 30.0 - 15.9 = \ang{14.1}$: una rotación en sentido horario del bloque
alrededor de un eje vertical. $P = \lambda_x - \lambda_o = -51.4 + 40.9 =
\ang{-10.5}$: la formación quedó \ang{10.5} más al norte de lo que indica la placa
estable, por lo que el bloque se ha desplazado desde entonces unos \ang{10.5}, unos
\qty{1170}{km}, hacia el sur respecto a la placa.
\end{solution}

\begin{solution}{exo:b3:palaeomagnetism:12}
(a) Con $\mathrm d\Omega = 2\pi\sin\theta'\,\mathrm d\theta'$,
\[
  \int_0^\theta\frac{\kappa e^{\kappa\cos\theta'}}{4\pi\sinh\kappa}\,2\pi\sin\theta'\,
  \mathrm d\theta' = \frac{e^\kappa - e^{\kappa\cos\theta}}{2\sinh\kappa}
  = \frac{e^\kappa - e^{\kappa\cos\theta}}{e^\kappa - e^{-\kappa}} .
\]
(b) Para $\kappa$ grande, $e^{-\kappa}$ es despreciable; dividiendo por $e^\kappa$
queda $1 - e^{-\kappa(1 - \cos\theta)} \approx 1 - e^{-\kappa\theta^2/2}$. Esto
vale $1 - e^{-1} = \qty{63}{\percent}$ cuando $\kappa\theta^2/2 = 1$, $\theta =
\sqrt{2/\kappa}$ radianes, o $\ang{81}/\sqrt\kappa$. (c) $81/\sqrt{50} =
\ang{11.5}$; la fórmula exacta da \qty{63.1}{\percent}.
\end{solution}

\begin{solution}{pb:b3:palaeomagnetism:1}
\textbf{1.} La hematites en capas rojas puede crecer a partir de óxidos e hidróxidos de hierro
mucho después de que se depositara el sedimento, dando lugar a una \omterm{def:b3:palaeomagnetism:nrm}{magnetización remanente química} de esa
edad posterior. Una prueba de pliegue positiva, si los estratos se plegaron poco después de
la sedimentación, o una prueba de inversión positiva entre estratos muy próximos de ambas
polaridades mostraría que la magnetización remanente es temprana.

\textbf{2.} $(\cos\ang{37}\cos\ang{50}, \cos\ang{37}\sin\ang{50}, \sin\ang{37})
= (0.5134, 0.6118, 0.6018)$.

\textbf{3.} $R = \sqrt{3.8943^2 + 2.6500^2 + 3.5717^2} = 5.9114$. $D =
\arctan(2.6500/3.8943) = \ang{34.2}$, $I = \arcsin(3.5717/5.9114) = \ang{37.2}$.

\textbf{4.} $k = 5/(6 - 5.9114) = 56.5$. Con $20^{1/5} - 1 = 0.8206$,
$\cos\alpha_{95} = 1 - (0.0886/5.9114)\times0.8206 = 0.98771$, $\alpha_{95} =
\ang{9.0}$.

\textbf{5.} B: $k = 6/(7 - 6.8979) = 58.8$; $20^{1/6} - 1 = 0.6475$,
$\cos\alpha_{95} = 1 - (0.1021/6.8979)\times0.6475 = 0.99042$, $\alpha_{95} =
\ang{7.9}$. C: $k = 7/(8 - 7.8707) = 54.1$; $20^{1/7} - 1 = 0.5340$,
$\cos\alpha_{95} = 1 - (0.1293/7.8707)\times0.5340 = 0.99123$, $\alpha_{95} =
\ang{7.6}$.

\textbf{6.} B tiene la $k$ mayor, el grupo más compacto; C tiene el
menor $\alpha_{95}$, la media mejor definida, porque tiene más localidades. Las tres
son similares: las tres formaciones están bien agrupadas.

\textbf{7.} $\lambda_A = \arctan(\tfrac12\tan\ang{37.17}) = \ang{20.76}$\,N,
$\lambda_B = \ang{32.15}$\,N, $\lambda_C = \ang{45.33}$\,N.

\textbf{8.} $p = 90 - 20.76 = \ang{69.24}$. $\sin\lambda_p =
\sin\ang{47}\cos\ang{69.24} + \cos\ang{47}\sin\ang{69.24}\cos\ang{34.23} =
0.2593 + 0.5272 = 0.7865$: $\lambda_p = \ang{51.86}$\,N.

\textbf{9.} $\sin\beta = \sin\ang{69.24}\sin\ang{34.23}/\cos\ang{51.86} =
0.8516$, $\beta = \ang{58.40}$. Dado que $\cos p = 0.3545 < \sin\ang{47}\times
0.7865 = 0.5752$, se aplica el segundo caso, $\cos\beta < 0$:
$\varphi_p = 15.00 + 180 - 58.40 = \ang{136.60}$\,E.

\textbf{10.} $\cos^2p = 0.1257$: $\mathrm dp = 4.50\times1.377 = \ang{6.2}$;
$\mathrm dm = 8.99\times\sin\ang{69.24}/\cos\ang{37.17} = \ang{10.5}$.

\textbf{11.} Desde \ang{51.9}\,N a \qty{60}{Ma} hasta \ang{66.5}\,N a
\qty{30}{Ma} y \ang{87.1}\,N a \qty{4}{Ma}, en casi el mismo meridiano
(\ang{136}--\ang{139}\,E): una trayectoria hacia el polo geográfico, alcanzado dentro
de su cono por la formación más joven.

\textbf{12.} La capa de ceniza, a \qty{22}{m}, se sitúa en la segunda zona normal,
19.7--\qty{24.6}{m}; con una edad cercana a \qty{3.2}{Ma}, esa zona es C2An.2n
(3.207--3.116). Luego, la zona 4.2--\qty{14.8}{m} es C2An.3n, la zona
29.2--\qty{54.0}{m} es C2An.1n; la base invertida pertenece al cron situado debajo de
C2An, y el techo invertido, al cron situado encima de él.

\textbf{13.} $(19.7 - 4.2)/(3.596 - 3.207) = 15.5/0.389 = \qty{39.8}{m/Myr}$;
$(54.0 - 24.6)/(3.116 - 2.581) = 29.4/0.535 = \qty{55.0}{m/Myr}$.

\textbf{14.} Base: $3.596 + 4.2/39.8 = \qty{3.70}{Ma}$. Techo: $2.581 -
10.0/55.0 = \qty{2.40}{Ma}$.

\textbf{15.} Solo se datan los límites de inversión, separados unos pocos cientos de miles de años
entre sí; entre ellos se supone constante la tasa. Un cambio en la subsidencia, en el
aporte de sedimento o en el clima, o una laguna en la sedimentación, la alteraría, y un
hiato podría incluso eliminar una zona corta y falsear el ajuste.

\textbf{16.} $\cos p = \sin\ang{47}\sin\ang{78} + \cos\ang{47}\cos\ang{78}
\cos\ang{175} = 0.7154 - 0.1413 = 0.5741$, $p = \ang{54.96}$. $\tan I_x =
2\cot p = 1.4023$, $I_x = \ang{54.5}$; $\lambda_x = \ang{35.04}$\,N.

\textbf{17.} $\cos D_x = (\sin\ang{78} - \sin\ang{47}\times0.5741)/
(\cos\ang{47}\sin\ang{54.96}) = 0.99976$; el polo está al este del
meridiano del yacimiento ($\sin\ang{175} > 0$), así que $D_x = \ang{1.3}$.

\textbf{18.} \omterm{def:b3:palaeomagnetism:vgp}{Paleolatitudes} esperadas a partir de $I_x$ = \ang{54.5}, \ang{60.0} y
\ang{63.6}: \ang{35.03}, \ang{40.89} y \ang{45.21}. $P_A = 35.03 - 20.76 =
\ang{14.3}$, $P_B = 40.89 - 32.15 = \ang{8.7}$, $P_C = 45.21 - 45.33 =
\ang{-0.1}$.

\textbf{19.} $R_A = 34.2 - 1.3 = \ang{32.9}$, $R_B = 23.9 - 2.9 = \ang{21.0}$,
$R_C = 3.4 - 2.5 = \ang{0.9}$, todas en sentido horario.

\textbf{20.} $\Delta D_o = \arcsin(\sin\ang{8.99}/\cos\ang{37.17}) = \ang{11.3}$;
$\Delta D_x = \arcsin(\sin\ang{3.0}/\cos\ang{54.5}) = \ang{5.2}$; $\Delta R_A =
\sqrt{11.3^2 + 5.2^2} = \ang{12.4}$.

\textbf{21.} $R_A = \ang{32.9} \pm \ang{12.4}$ y $R_B = \ang{21.0} \pm
\ang{13.8}$ son significativos; $R_C = \ang{0.9} \pm \ang{17.9}$ no lo es.

\textbf{22.} Hace sesenta millones de años, Kesta se encontraba unos \ang{14} de latitud
al sur de donde lo sitúa el continente estable, y estaba girado unos \ang{33}
en sentido antihorario respecto a su orientación actual. Para \qty{30}{Ma} había avanzado
\ang{6} de esos \ang{14} hacia el norte y girado \ang{12} de esos \ang{33} en sentido horario;
para \qty{4}{Ma} había llegado a su lugar actual en el continente,
sin rotación pendiente. Derivó hacia el norte a lo largo del margen del continente, girando
en sentido horario a medida que avanzaba, y ya estaba en su lugar hacia \qty{4}{Ma}.

\textbf{23.} $47.00 - 20.76 = \ang{26.24}$, o $26.24\times111.19 =
\qty{2918}{km}$. Esto incluye el movimiento del propio continente estable;
respecto a él, Kesta se desplazó $P_A = \ang{14.3}$.

\textbf{24.} $I = \ang{28.18}$ y \ang{46.16} dan $\lambda = \ang{15.00}$ y
\ang{27.50}: una semianchura de \ang{6.25}, y una desviación típica de
$6.25/1.96 = \ang{3.19}$, o \qty{355}{km}. Desde la formación A, Kesta ha derivado
hacia el norte en
\textbf{\qty[text-series-to-math=true]{2918 \pm 355}{km}}.
\end{solution}
